\title{xLSTM: Extended Long Short-Term Memory}

\begin{abstract}
The introduction of the constant error carousel and gating mechanisms in the 1990s established the foundational principles behind Long Short-Term Memory (LSTM) networks. These innovations allowed LSTMs to remain highly relevant within the deep learning landscape, notably facilitating the development of the earliest Large Language Models (LLMs). Nevertheless, the rise of Transformer architectures—characterized by inherently parallelizable self-attention—signaled a significant technological shift, eventually surpassing LSTMs when scaling to large models. This work seeks to explore a fundamental question: By scaling LSTMs up to the order of billions of parameters and integrating advances from modern LLMs, while also addressing existing limitations inherent to LSTMs, to what extent can LSTMs compete in the domain of language modeling? To this end, we present exponential gating, accompanied by suitable normalization and stabilization methods. Moreover, we introduce modifications to the LSTM’s memory construct, resulting in the following architectures: (i) sLSTM, incorporating scalar memory and updates as well as novel memory mixing, and (ii) mLSTM, which permits full parallelization using a matrix memory paired with a covariance-based update rule. Embedding these enhanced LSTM variants within residual block structures leads to the formation of xLSTM blocks, which are then assembled into residual-stacked xLSTM networks. These innovations—exponential gating and revised memory configurations—enable xLSTM models to rival state-of-the-art Transformer and State Space Model approaches, excelling both in performance and scalability.\\
Code available at: \url{https://github.com/NX-AI/xlstm}
\end{abstract}

\section{Introduction}
\label{sec:intro}

Long Short-Term Memory (LSTM) concepts~\citep{Hochreiter:91,Hochreiter:97nips,Hochreiter:97}, specifically the constant error carousel and gating mechanisms, were developed as a solution to the vanishing gradient issue in recurrent neural networks~\citep{Hochreiter:91,Hochreiter:00book}:

\begin{align}
\label{eq:lstm_idea}
  c_t \ = \ f_t \ c_{t-1} \ + \ 
          i_t \  z_t \ , \quad  
          h_t \ = \ o_t \ \psi({c_t}) \ . 
\end{align}

The constant error carousel involves incrementally modifying the previous cell state $c_{t-1}$ (green) with inputs $z_t$, with this process regulated by sigmoid gates (blue). The input gate $i_t$ and the forget gate $f_t$ jointly determine how the cell state is updated, whereas the output gate $o_t$ manages the cell’s output, specifically the hidden state $h_t$. Before producing the hidden state, the cell state undergoes normalization or squashing via $\psi$, and the outcome is filtered by the output gate.

LSTMs have found wide-ranging applications across numerous fields \citep{Hochreiter:01,Hochreiter:07,Schmidhuber:15}, maintaining dominance in text generation tasks up to the emergence of Transformers in 2017~\citep{Vaswani:17}. Their utility has been substantiated in various sequence modeling challenges, including but not limited to text generation~\citep{Graves:13,Karpathy:15blog}, handwriting synthesis~\citep{Graves:13}, sequence-to-sequence translation~\citep{Sutskever:14nips}, computer program evaluation~\citep{Zaremba:14arxiv}, image captioning~\citep{Karpathy:15,Hossain:19}, code generation~\citep{Karpathy:15blog}, rainfall-runoff forecasting~\citep{Kratzert:18,Kratzert:19}, and the development of hydrological flood-warning systems~\citep{Nearing:24}. Within the realm of reinforcement learning, LSTMs remain the sequence models of choice, powering leading systems such as AlphaStar for StarCraft II~\citep{Vinyals:17short}, OpenAI Five for Dota 2~\citep{OpenAI:19}, and magnetic controllers for nuclear fusion~\citep{Degrave:22short}. The strength of LSTMs lies in their capacity for abstraction—skilfully distilling and retaining semantic content within memory cells~\citep{Karpathy:15blog}. This ability has been evidenced by the discovery of specialized neurons responsible for tasks such as recognizing numbers and syntax~\citep{Lakretz:19}, linguistic features~\citep{Bau:19}, and sentiment~\citep{Radford:17arxiv}. LSTMs continue to play a crucial role in contemporary and impactful applications~\citep{Degrave:22short,Nearing:24}, demonstrating their enduring relevance in sequence modeling.

Although LSTMs have achieved remarkable results, they are constrained by three significant shortcomings: (i) They are unable to revise previously made storage choices. This is illustrated through the {\it Nearest Neighbor Search} task (see also Appendix~\ref{sec:appNearestNeighborSearch}), where, given a reference vector, the sequence must be iterated through step by step to locate the vector most similar to the reference, returning the value associated with it at the end. The left side of Figure~\ref{fig:lstmProblems} displays the mean squared error for this task. LSTM models find it difficult to update a stored value if a more closely matching vector appears later in the sequence, in contrast to our proposed xLSTM, which addresses this issue through the use of exponential gating mechanisms. (ii) They suffer from constrained storage capacity because all information is compressed into single scalar cell states. This limitation is evident in the {\it Rare Token Prediction} task. As depicted in the right side of Figure~\ref{fig:lstmProblems}, token prediction perplexity on Wikitext-103~\citep{Merity:17} is reported for different frequency partitions. LSTMs exhibit inferior performance on infrequent tokens, an issue attributable to their limited memory resources. Our xLSTM model overcomes this challenge by employing matrix-based memory representations. (iii) The architecture inherently lacks parallel processing ability owing to memory mixing—hidden-to-hidden dependencies across time steps force strictly sequential computation.

The shortcomings inherent in LSTM architectures have catalyzed the development and adoption of Transformers~\citep{Vaswani:17} for language modeling tasks. This raises the question: by addressing these issues and expanding LSTM models to match the scale of today's Large Language Models, what levels of performance might be attainable in language modeling?

\section{Extended Long Short-Term Memory}
\label{sec:xLSTM}

To address the limitations inherent in LSTM architectures, Extended Long Short-Term Memory (xLSTM) incorporates two significant innovations into the foundational concept presented in Equation~\eqref{eq:lstm_idea}. Specifically, xLSTM introduces exponential gating alongside novel memory architectures, thereby expanding the LSTM framework with two new variants: (i) the sLSTM (detailed in Section~\ref{sec:sLSTM}), characterized by a scalar memory, a scalar update process, and the presence of memory mixing, and (ii) the mLSTM (see Section~\ref{sec:mLSTM}), which utilizes a matrix-based memory coupled with an update mechanism relying on covariance (via the outer product), resulting in an approach that can be entirely parallelized. Both of these models extend LSTM functionality through the use of exponential gating. In order to facilitate parallelization, memory mixing (i.e., hidden-to-hidden recurrence) is omitted in the mLSTM architecture. The scalability of both sLSTM and mLSTM to configurations with multiple memory cells is possible, with the sLSTM supporting inter-cell memory mixing. Additionally, the sLSTM allows for multiple heads without cross-head memory mixing, although memory mixing can occur among cells within a single head. The introduction of this multi-head structure in sLSTM, in conjunction with exponential gating, provides a novel mechanism for orchestrating memory interactions. For the mLSTM, the notions of multiple heads and multiple cells are functionally interchangeable.

By embedding these novel LSTM variants within residual block structures, we construct xLSTM blocks (refer to Section~\ref{sec:xLSTMblock}). Building architectures by residually layering multiple xLSTM blocks leads to the creation of xLSTM architectures (see Section~\ref{sec:xLSTMarchitecure}).
The xLSTM architecture along with its constituent components is illustrated in Figure~\ref{fig:xlstm_sketch}.

\subsection{Review of the Long Short-Term Memory}
\label{sec:orgLSTM}

The initial LSTM framework~\citep{Hochreiter:91,Hochreiter:97nips,Hochreiter:97} conceptualized the memory cell as a core unit responsible for both computation and storage, effectively mitigating the vanishing gradient problem~\citep{Hochreiter:91,Hochreiter:00book} via the mechanism known as the constant error carousel (cell state update). This architecture incorporates three distinct gates—input, output, and forget gates—with the latter being introduced later by \citet{Gers:00}. The LSTM cell’s dynamics at a given time step $t$ are governed by the following update equations:

\begin{align}
c_t \ &= \  f_t \ c_{t-1} \ + \ i_t \ z_t &  & &\text{cell state} \\
h_t  \ &= \ o_t \ \tilde{h}_t \ , & \tilde{h}_t \ &= \ \psi \left( c_t\right)
  &\text{hidden state} \\
z_t \ &= \ \varphi \left( \tilde{z}_t \right) \ , 
  &\tilde{z}_t \ &=  \ \mathbf{w}^\top_{z} \ \mathbf{x}_t \ + \
  r_{z}  h_{t-1} \ + \  b_{z} \ \
  &\text{cell input} \\
i_t \ &= \ \sigma \left( \tilde{i}_t  \right) \ , 
  &\tilde{i}_t \ &= \ \mathbf{w}^\top_{i} \ \mathbf{x}_t \ + \
  r_{i}  \ h_{t-1} \ + \  b_{i} \ \
  &\text{input gate} \\
f_t \ &= \ \sigma \left(  \tilde{f}_t \right) \ , 
  &\tilde{f}_t \ &= \ \mathbf{w}^\top_{f} \ \mathbf{x}_t  \ + \
  r_{f}  \ h_{t-1} \ + \  b_{f} \ \
  &\text{forget gate} \\
o_t \ &= \ \sigma \left( \tilde{o}_t \right) \ , 
  &\tilde{o}_t  \ &= \ \mathbf{w}^\top_{o} \ \mathbf{x}_t \ + \
  r_{o}  \ h_{t-1} \ + \  b_{o} \ \
  &\text{output gate} 
\end{align}

The vectors $\mathbf{w}_{z}$, $\mathbf{w}_{i}$, $\mathbf{w}_{f}$, and $\mathbf{w}_{o}$ denote the sets of weights connecting the inputs $\mathbf{x}_t$ to the cell input, input gate, forget gate, and output gate, respectively. The parameters $r_{z}$, $r_{i}$, $r_{f}$, and $r_{o}$ represent the recurrent weight connections from the previous hidden state $h_{t-1}$ to each of these components. Biases associated with each unit are indicated by $b_{z}$, $b_{i}$, $b_{f}$, and $b_{o}$. The functions $\varphi$ and $\psi$ correspond to the activations for the cell input and hidden state, most often chosen as $\tanh$. The purpose of $\psi$ is to constrain the cell state, which can grow without bound if left unchecked. All gating mechanisms implement a sigmoid activation function, $\sigma \left( x \right)= 1/(1+\exp(-x))$. 

In more recent developments, multiple memory cell values are assembled into a vector $\mathbf{c}_t \in \mathbb{R}^{d}$, enabling the use of recurrent weight matrices $\mathbf{R} \in \mathbb{R}^{d \times d}$ to combine or mix the outputs of the different memory cells~\citep{Greff:15}; for further explanation, see Appendix~\ref{sec:apporgLSTM}. Comprehensive ablation experiments have demonstrated that every constituent of the memory cell plays an essential role~\citep{Greff:15}.

\subsection{sLSTM}
\label{sec:sLSTM}

In order to enable LSTMs to modify their storage choices dynamically, we incorporate exponential gates (depicted in red), complemented by normalization and stabilization mechanisms. Specifically, the activation functions of both the input and forget gates are allowed to be exponential. To implement normalization, we define a normalizer state that accumulates the products of the input gate and each subsequent forget gate.

The scalar sLSTM forward pass is:

\begin{align}
\label{eq:slstmforward}
c_t \ &= \  f_t \ c_{t-1} \ + \ i_t \ z_t & & 
 &\text{cell state} \\
n_t \ &= \  f_t \ n_{t-1} \ + \ i_t  & & 
 &\text{normalizer state} \\
h_t \ &= \ o_t \ \tilde{h}_t \ ,
  &\tilde{h}_t \ &= \ c_t / n_t
&\text{hidden state} \\
z_t \ &= \ \varphi \left( \tilde{z}_t \right) \ , 
  &\tilde{z}_t \ &=  \ \mathbf{w}^\top_{z} \ \mathbf{x}_t \ + \
  r_{z}  h_{t-1} \ + \  b_{z}
  &\text{cell input} \\
i_t \ &= \ \!\exp \left( \tilde{i}_t  \right) \ , 
  &\tilde{i}_t \ &= \ \mathbf{w}^\top_{i} \ \mathbf{x}_t \ + \
  r_{i}  \ h_{t-1} \ + \  b_{i}
  &\text{input gate} \\
f_t \ &= \ \sigma \left(  \tilde{f}_t \right) \ \text{OR} \
\!\exp \left(  \tilde{f}_t \right) \ ,
  &\tilde{f}_t \ &= \ \mathbf{w}^\top_{f} \ \mathbf{x}_t  \ + \
  r_{f}  \ h_{t-1} \ + \  b_{f}
  &\text{forget gate} \\
o_t \ &= \ \sigma \left( \tilde{o}_t \right) \ , 
  &\tilde{o}_t  \ &= \ \mathbf{w}^\top_{o} \ \mathbf{x}_t \ + \
  r_{o}  \ h_{t-1} \ + \  b_{o}
  &\text{output gate} 
\end{align}

We adapt the classic LSTM gating mechanisms—specifically, input- and/or hidden-driven gates along with bias parameters—into these novel network designs. As exponential activations may generate extremely large outputs and potential overflows, we introduce a stabilizing term for the gates via an auxiliary state \mbox{$m_t$~\citep{Milakov:18arxiv}}:

\begin{align}
\label{eq:slstmstabil}
m_t \ &= \ \max \left( \log ( f_t ) + m_{t-1} , \log ( i_t ) \right) &\text{stabilizer state} \\
i'_t \ &= \ \!\exp \left( \log \left ( i_t \right) - m_t \right) \ = \ \!\exp \left( \tilde{i}_t - m_t \right) \ \, 
  &\text{stabil. input gate} \\
f'_t \ &= \ \!\exp \left( \log \left( f_t \right) + m_{t-1} - m_t \right) \ \, 
  &\text{stabil. forget gate}
\end{align}

As demonstrated in Appendix~\ref{sec:appsLSTM}, substituting $f_t$ with $f'_t$ and $i_t$ with $i'_t$ during the forward propagation leaves both the network's overall output and the gradients of the loss with respect to the parameters unchanged.

\paragraph{New Memory Mixing.} Similar to the original LSTM, sLSTM supports the use of multiple memory cells (refer to Appendix~\ref{sec:appsLSTM}). The presence of several memory cells allows for the mixing of memory through recurrent pathways, specifically $\mathbf{R}_{\mathbf{z}}$, $\mathbf{R}_{\mathbf{i}}$, $\mathbf{R}_{\mathbf{f}}$, $\mathbf{R}_{\mathbf{o}}$, which connect the hidden state vector $\mathbf{h}$ to the memory cell input $\mathbf{z}$ as well as the gating vectors $\mathbf{i}$, $\mathbf{f}$, and $\mathbf{o}$. What distinguishes this approach to memory mixing is the integration of exponential gating mechanisms. In the updated sLSTM architecture, multiple heads are permitted, where memory mixing occurs within each head independently but not between heads. By combining these heads with exponential gating, sLSTM introduces a novel strategy for memory mixing.

\subsection{mLSTM}
\label{sec:mLSTM}

To improve the storage potential of LSTM architectures, we substitute the traditional scalar memory cell $c \in \mathbb{R}$ with a matrix memory cell $\mathbf{C} \in \mathbb{R}^{d \times d}$. As a result, memory access takes place through matrix-matrix operations. At each time step $t$, the system must store a pair of vectors: the key $\mathbf{k}_t \in \mathbb{R}^d$ and the value $\mathbf{v}_t \in \mathbb{R}^d$, following the nomenclature commonly used in Transformer models. Subsequently, at a later time $t + \tau$, a query vector $\mathbf{q}_{t+\tau} \in \mathbb{R}^d$ should enable the retrieval of the previously stored value $\mathbf{v}_t$. This framework aligns with the classic scenario of Bidirectional Associative Memories (BAMs) \citep{Kohonen:72,Anderson:72,Nakano:72,Anderson:77}.

The covariance update rule~\citep{Sejnowski:77,Dayan:91} for storing a key-value pair is

\begin{align}
\mathbf{C}_t \ &= \ \mathbf{C}_{t-1} \ + \ \mathbf{v}_{t} \ \mathbf{k}_t^\top \ .
\end{align}

We posit that a layer normalization is applied before projecting the inputs into key and value representations, ensuring their means are zero. The update rule for the covariance is theoretically optimal~\citep{Dayan:91} for achieving maximal separability of recovered binary vectors; this goal aligns with maximizing the signal-to-noise ratio. When considering only pairwise retrievals and permitting the quadratic computational cost inherent to attention mechanisms~\citep{Krotov:16,Krotov:17,Ramsauer:21}, separability can be further increased. The covariance update rule shares its formulation with Fast Weight Programmers \citep{Schmidhuber:92ncfastweights,Schlag:21}, where subsequent approaches have introduced a constant decay coefficient multiplying $\mathbf{C}_{t-1}$ and a constant learning coefficient multiplying 
$\mathbf{v}_{t} \mathbf{k}_t ^\top$~\citep{Ba:16}.
Following this line of reasoning, we adapt the covariance update procedure within the LSTM architecture: the forget gate operates as the decay parameter, the input gate modulates the learning rate, and the output gate governs the scale of the returned vector.

In this matrix memory setup, the normalizer state is constructed as a weighted aggregation of key vectors, with weights determined by both the input gate and the multiplicative effects of all succeeding forget gates. This normalizer state effectively tracks the cumulative influence of the gating mechanisms. Because the dot product between the query vector and the normalizer state can approach zero, we instead employ the magnitude of this dot product, enforcing a lower bound (commonly set to 1.0) to ensure stability, following the approach in~\citep{Sun:23arxiv}. The resulting mLSTM forward computation proceeds as follows:

\begin{align}
\label{eq:mlstm_recurrent_begin}
\mathbf{C}_t \ &= \  f_t \ \mathbf{C}_{t-1} \ + \ 
  i_t \ \mathbf{v}_t \ \mathbf{k}_t^\top & &  &\text{cell state} \\
\mathbf{n}_t \ &= \  f_t \ \mathbf{n}_{t-1} \ + \ 
  i_t \ \mathbf{k}_t & &  &\text{normalizer state} \\
\mathbf{h}_t  \ &= \ \mathbf{o}_t \ \odot \ \tilde{\mathbf{h}}_t
\ , \qquad \qquad 
\tilde{\mathbf{h}}_t \ = \   \mathbf{C}_t \mathbf{q}_t \ / \ 
  \max \left\{ |\mathbf{n}_t^\top \mathbf{q}_t|, 1 \right\} 
   &&&\text{hidden state} \\
\mathbf{q}_t \ &= \ \mathbf{W}_q \ \mathbf{x}_t \ + \ \mathbf{b}_q  & &  &\text{query input} \\
\mathbf{k}_t \ &= \ \frac{1}{\sqrt{d}} \mathbf{W}_k \ \mathbf{x}_t \ + \ \mathbf{b}_k  & &  &\text{key input} \\
\mathbf{v}_t \ &= \ \mathbf{W}_v \ \mathbf{x}_t \ + \ \mathbf{b}_v  & &  &\text{value input} \\
i_t \ &= \ \!\exp \!  \! \left( \tilde{i}_t  \right) \ , 
 \qquad \qquad \ \ \ \ \,
  \tilde{i}_t \ = \ \mathbf{w}^\top_{i} \ \mathbf{x}_t \ + \  b_{i} \ \ \ 
  &&&\text{input gate} \\
f_t \ &= \ \sigma \!  \left(  \tilde{f}_t \right) \ \text{OR} \ \!\exp \!  \! \left(  \tilde{f}_t \right) , \ \ \: \!
\tilde{f}_t \ = \ \mathbf{w}^\top_{f} \ \mathbf{x}_t  \ + \
  b_{f}
  &&& \text{forget gate} \\
\label{eq:mlstm_recurrent_end}
\mathbf{o}_t \ &= \ \sigma \left( \tilde{\mathbf{o}}_t \right) \ , \qquad \qquad \quad \ \ \ \,
  \tilde{\mathbf{o}}_t  \ = \ \mathbf{W}_{\mathbf{o}} \ \mathbf{x}_t \ + \
  \mathbf{b}_{\mathbf{o}}
  &&&\text{output gate} 
\end{align}

Similar to the conventional LSTM, mLSTM supports several memory cells. In the case of mLSTM, having multiple heads functions equivalently to having multiple memory cells, due to the absence of memory mixing. To ensure the stability of the exponential gating mechanisms within mLSTM, we adopt the stabilization methods employed in sLSTM (refer to Equation~\ref{eq:slstmstabil}). Owing to the lack of memory mixing in mLSTM, its recurrence relations can be expressed in a parallelizable form. Additional information can be found in Appendix~\ref{sec:appmLSTM}.

\subsection{xLSTM Architecture}
\label{sec:xLSTMarch}

\paragraph{xLSTM Blocks.}
\label{sec:xLSTMblock}
An xLSTM block aims to capture and non-linearly encode the sequence history within a high-dimensional representation, enhancing its capacity to disentangle distinct past contexts or histories. Distinguishing between such histories is essential for accurately forecasting subsequent elements in a sequence, such as predicting the next token. Our approach draws on Cover’s Theorem~\citep{Cover:65}, which posits that patterns embedded through nonlinear mappings in high-dimensional spaces are more likely to become linearly separable compared to their configurations in lower dimensions. 

We analyze two types of residual block designs: (i) A residual block featuring a post up-projection architecture (similar to that used in Transformers), where the sequence history is first summarized non-linearly within its original space, then subjected to a linear up-projection into a high-dimensional space, followed by a nonlinear activation and subsequently a linear projection returning it to the original space. This is illustrated in the leftmost panel of Figure~\ref{fig:Backbones} and the third column of Figure~\ref{fig:xlstm_sketch}. A more comprehensive schematic is available in Figure~\ref{fig:appsLSTM_detailed} in the appendix. (ii) The alternative is a residual block using pre up-projection (analogous to State Space Models), which initiates with a linear projection into a higher-dimensional space,

It performs a nonlinear aggregation of historical information within a high-dimensional space, followed by a linear projection back to the input space. When an xLSTM block utilizes an sLSTM, we apply the up-projection after the block. Conversely, for xLSTM blocks using an mLSTM, the up-projection is applied beforehand, taking advantage of the increased memory capacity afforded by the high-dimensional representation. For further clarification, see the left section of Figure~\ref{fig:Backbones}, the third column of Figure~\ref{fig:xlstm_sketch}, or consult Figure~\ref{fig:appsLSTM_detailed} in the appendix.

\paragraph{xLSTM Architecture. }
\label{sec:xLSTMarchitecure}
The xLSTM model is formed by stacking individual modules in a residual manner~\citep{Srivastava:15,He:16}. This design adopts the widely used pre-LayerNorm~\citep{Ba:16b} residual backbone, which is the standard in current Large Language Models. For illustration, refer to the final column of Figure~\ref{fig:xlstm_sketch}.

\subsection{Memory and Speed Considerations}

In contrast to Transformers, xLSTM networks feature linear computational complexity and maintain fixed memory requirements regardless of sequence length. Owing to the compressive nature of xLSTM memory, these networks are particularly advantageous for industrial use cases and edge deployment.

Unlike parameter-dependent memory architectures, the memory in mLSTM operates without needing parameters, but managing and updating its $d \times d$ matrix memory comes at a high computational cost. Our approach negotiates this trade-off by balancing memory size with computational demands. Importantly, these operations are well-suited to parallel processing on GPUs, resulting in a minimal impact on actual run time.

Although mLSTM supports parallelization similar to FlashAttention~\citep{Dao:22,Dao:23} and GLA~\citep{Yang:23arxiv}, the sLSTM architecture is inherently sequential as a result of its memory mixing via hidden-to-hidden connections, preventing parallel execution. Nevertheless, we introduced an efficient CUDA implementation incorporating register-level GPU memory optimizations, resulting in execution speeds typically no more than twice as slow as those of mLSTM.

\section{Related Work}
\label{sec:related}

\paragraph{Linear Attention.} To address the quadratic computational complexity of the Transformer with respect to context length, a range of approaches have been developed to achieve linear scalability. The Synthesizer framework generates learned, synthetic attention scores without modeling explicit token-to-token dependencies~\citep{Tay:20arxiv}. Linformer reduces the dimensionality of self-attention through a low-rank projection, thereby enabling a linear approximation~\citep{Wang:20arxiv}. In a similar direction, the Linear Transformer modifies the attention mechanism to achieve linear computation~\citep{Katharopoulos:20}. Performer introduces a linear approximation of the attention softmax by employing positive orthogonal random features~\citep{Choromanski:21}. Additionally, conventional attention has been substituted by rapid long-sequence convolutions in architectures such as the Structured Global Convolution (SGConv)~\citep{Li:22} and the Hyena Hierarchy~\citep{Poli:23}.

\paragraph{State Space Models.} In recent years, State Space Models (SSMs) have gained considerable attention due to their linear complexity with respect to sequence length and their competitive results against Transformer architectures. Early examples include the Structured State Space sequence model (S4)~\citep{Gu:21}, which was subsequently followed by models such as the Diagonal State Space (DSS)~\citep{Gupta:22}, Gated State Space (GSS)~\citep{Mehta:22}, S5~\citep{Smith:22}, Bidirectional Gated SSM (BiGS)~\citep{Wang:22}, H3~\citep{Fu:23}, and Mamba~\citep{Gu:24arxiv}.

\paragraph{Recurrent Neural Networks.} Recurrent Neural Networks (RNNs) have been proposed as alternatives to Transformer architectures and attention mechanisms, primarily because they scale linearly with the context length. Deep Linear Recurrent Units (LRUs) within RNNs have yielded strong performance on language modeling tasks~\citep{Orvieto:23, De:24arxiv}, as have other RNN variants such as Hierarchically Gated Linear RNN (HGRN)~\citep{Qin:23} and its successor HGRN2~\citep{Qin:24arxiv}. Among prominent RNN models for large-scale language modeling is RWKV~\citep{Peng:23arxivshort,Peng:24arxivshort}, which has demonstrated results comparable to Transformer-based models.

\paragraph{Gating.}
Gating, a fundamental concept originally introduced in LSTM architectures, has subsequently been revived and adapted across various modern methods. Its implementation appears in numerous recent models, such as HGRN~\citep{Qin:23}, HGRN2~\citep{Qin:24arxiv}, Gated Linear Attention (GLA)~\citep{Yang:23arxiv}, Gated State Space (GSS) models~\citep{Mehta:22}, Bidirectional Gated SSM (BiGS)~\citep{Wang:22}, Moving Average Equipped Gated Attention (MEGA)~\citep{Ma:22}, RWKV~\citep{Peng:23arxivshort}, and Mamba~\citep{Gu:24arxiv}.

\paragraph{Covariance Update Rule.} In order to improve storage efficiency, we integrated a matrix memory into the mLSTM cell that employs a covariance-based update mechanism. This approach has also been adopted in several related works, including Fast Weight Programmers \citep{Schmidhuber:92ncfastweights,Schlag:21}, RWKV-5 and RWKV-6~\citep{Peng:24arxivshort}, Retention~\citep{Sun:23arxiv}, Linear Transformer~\citep{Katharopoulos:20}, and HGRN2~\citep{Qin:24arxiv}.

\paragraph{Most Related.} Among existing architectures, Retention~\citep{Sun:23arxiv}, RWKV~\citep{Peng:23arxivshort,Peng:24arxivshort}, and HGRN2~\citep{Qin:24arxiv} are most closely related to xLSTM, as they employ mechanisms such as matrix memory and gating. Nevertheless, unlike the proposed sLSTM, these models do not implement memory mixing. The incorporation of memory mixing is crucial for effectively addressing state tracking tasks, giving LSTMs a greater expressive capacity compared to State Space Models (SSMs) and Transformers~\citep{Merrill:24,Deletang:23}. Accurately tracking state is particularly essential when evaluating code or following entities across lengthy narratives.

\paragraph{Residually Stacking Architectures.} Similar to nearly all modern, large-scale deep learning models, the xLSTM frameworks are developed by stacking their core modules with residual connections~\citep{Srivastava:15,He:16}.

This design approach was foundational in enabling deep convolutional architectures~\citep{He:16} as well as the emergence of the Transformer~\citep{Vaswani:17}. In turn, the Transformer serves as the essential foundation for prominent Large Language Models (LLMs) such as GPT-3~\citep{Brown:20short}, ChatGPT~\citep{Schulman:22short}, GPT-4~\citep{Achiam:23gpt4short}, Megatron-LM~\citep{Shoeybi:19}, Gopher~\citep{Rae:21short}, ERNIE 3.0 Titan~\citep{Wang:21arxivshort}, GLaM~\citep{Du:21glamshort}, Chinese M6~\citep{Lin:21}, the multilingual AlexaTM 20B~\citep{Soltan:22}, OPT~\citep{Zhang:22opt}, Chinchilla~\citep{Hoffmann:22short}, BLOOM~\citep{Scao:22arxivshort}, GLM-130B~\citep{Zeng:22arxivshort}, LaMDA~\citep{Thoppilan:22short}, PaLM~\citep{Chowdhery:22short}, Llama~\citep{Touvron:23arxiv}, and Gemini~\citep{GeminiTeam:23arxiv,Reid:24arxiv}.

\section{Experiments}
\label{sec:Exp}

This section presents an empirical analysis of xLSTM, benchmarking its performance against established approaches, with particular attention to language modeling tasks. Section~\ref{sec:ExpSyntethic} explores the distinct capabilities of xLSTM by applying it to synthetic benchmarks. 

Next, Section~\ref{sec:ExpComparison} details a comparison of language modeling techniques based on their validation set perplexity, using models trained on a subset of 15B tokens from the SlimPajama dataset~\citep{Soboleva:23}. Within this framework, we also conduct ablation studies focused on the xLSTM architecture. Additionally, this section investigates the scaling properties of all evaluated models following methodologies similar to those in \citet{Kaplan:20} and \citet{Brown:20short}. 

Further, in Section~\ref{sec:ExpLanguage}, we extend our language modeling investigation by training both xLSTM and the top-performing architectures identified previously on an expanded dataset comprising 300B SlimPajama tokens~\citep{Soboleva:23}. The evaluation procedure encompasses: (1) measuring each model's generalization to longer context lengths; (2) analyzing validation perplexity and performance across a suite of downstream tasks~\citep{Sutawika:24}; (3) testing on 571 textual domains from the PALOMA language benchmark~\citep{Magnusson:23arxivshort}; and (4) reassessing scaling behaviors with this significantly larger training set, now using 20-fold more data.

Throughout our experiments, we adopt the notation xLSTM[$a$:$b$] to indicate the proportion $a/b$ between mLSTM-based and sLSTM-based xLSTM blocks. As an illustration, xLSTM[7:1] denotes a composition where, among eight total blocks, seven are implemented as mLSTM-based blocks, and one utilizes an sLSTM-based architecture. When scaling to a standard configuration of 48 total blocks, this notation corresponds to allocating 42 mLSTM-based blocks and 6 sLSTM-based blocks. In addition, in every experiment, we consistently incorporate up-projection blocks prior to mLSTM and following sLSTM components.

\subsection{Synthetic Tasks and Long Range Arena}
\label{sec:ExpSyntethic}
We begin by evaluating the novel exponential gating with memory mixing mechanism of xLSTM on formal languages~\citep{Deletang:23}. Next, we investigate the capabilities of xLSTM’s innovative matrix memory component through the Multi-Query Associative Recall task~\citep{Arora:23arxiv}. Lastly, the capacity of xLSTM to handle long sequence data is measured using benchmarks from the Long Range Arena~\citep{Tay:21}.

\paragraph{Evaluation of xLSTM's Exponential Gating and Memory Mixing.} We assess the efficacy of the exponential gating with memory mixing introduced in xLSTM, a mechanism designed to support state tracking challenges~\citep{Merrill:24, Merrill:23}. Building upon the suite of formal language benchmarks from \citet{Deletang:23}, we extend these tasks to incorporate training across sequences of varying lengths to facilitate length extrapolation. A comprehensive account of the experimental setup and supplemental findings can be found in Appendix~\ref{sec:appExpSynthetic-formalLang}. xLSTM is benchmarked against a range of alternative models, including Transformers, State Space Models, and standard Recurrent Neural Networks. Performance is assessed by measuring the accuracy on tokens pertinent to the target task, using a scoring scale from 0 (indicating chance) to 1 (denoting flawless performance). We examine 2-block configurations for several models on these benchmarks: xLSTM[0:1] (pure sLSTM), xLSTM[1:0] (pure mLSTM), xLSTM[1:1], Llama, Mamba, RWKV, Retention, Hyena, LSTM, and a hybrid approach with LSTM components in Transformer blocks (LSTM (Block)). Experimental outcomes are displayed in Figure~\ref{fig:formal-main}. Models such as Transformers or State Space Models that lack memory mixing (and therefore state tracking ability) are, for example, unable to handle regular grammar tasks like parity. This is consistent with previous work that demonstrates the fundamental limitations of Transformers and State Space Models in comparison to RNNs~\citep{Merrill:24, Merrill:23, Deletang:23}.

\paragraph{Test of xLSTM's Memory Capacities on Associative Recall Tasks.} In this study, we evaluate the matrix-based memory mechanism of xLSTM by assessing its ability to handle the Multi-Query Associative Recall task~\citep{Arora:23arxiv}: Each input sequence comprises randomly selected key-value pairs from an extensive vocabulary that the model is required to retain for subsequent recall. To further challenge the original task, we scale the number of key-value pairs up to 256 and extend the sequence context length to 2048, enabling a more comprehensive assessment of each model's memory limits. We analyze the performance of 2-block versions of Llama, Mamba, RWKV-5, RWKV-6, xLSTM[1:1], and xLSTM[1:0] by measuring their recall accuracy. Transformers such as Llama, which exhibit memory scaling exponentially with respect to the coding dimension~\citep{Ramsauer:21}, serve as a performance benchmark for this evaluation. The comparative outcomes are depicted in Figure~\ref{fig:mqar-main}. 
Among models outside the Transformer family, xLSTM[1:1] consistently achieves superior results, including for models with smaller capacities. Notably, the inclusion of the sLSTM block does not impede recall capacity; instead, it appears to improve performance, particularly apparent in the most complex scenario involving 256 key-value pairs. Appendix~\ref{sec:appExpSynthetic-mqar} provides further findings, with extrapolation studies suggesting that xLSTM's improved memory capacities enable generalization to longer context lengths than those used in training.

\paragraph{Test of xLSTM's Long Context Capabilities on Long Range Arena.} To evaluate how well xLSTM manages extended sequences and broad contextual information, we conduct a comparison with alternative approaches using the Long Range Arena~\citep{Tay:21}. Across all tasks, xLSTM achieves reliably high results, indicating that its architecture is particularly adept at addressing a wide range of long-context challenges.

For more details, see Appendix~\ref{sec:appExpSynthetic-lra}.

\subsection{Method Comparison and Ablation Study}
\label{sec:ExpComparison}

In order to explore the central question posed in this study—specifically, the potential performance of our novel LSTM variants when applied to large-scale language modelling—we conduct experiments where xLSTMs, Transformers, State Space Models, and additional approaches are all trained on a corpus of 15B tokens drawn from SlimPajama, utilizing a consistent auto-regressive framework. Performance evaluations are carried out on a validation dataset, and further ablation analyses are undertaken specifically for the xLSTM models.

\paragraph{Comparing xLSTM to Other Methods.} To provide a comprehensive comparison, we train each model using 15B tokens from the SlimPajama dataset~\citep{Soboleva:23}. Evaluation is conducted by measuring perplexity on a held-out validation set. The following architectures are considered: our proposed xLSTM, GPT-3 (Transformer-based)~\citep{Brown:20short}, Llama (Transformer-based)~\citep{Touvron:23arxiv}, H3 (SSM category)~\citep{Fu:23}, Mamba (SSM category)~\citep{Gu:24arxiv}, RWKV-4, RWKV-5, and RWKV-6 (all RNN-based)~\citep{Peng:23arxivshort,Peng:24arxivshort}, GLA (a variant of linear Transformer)~\citep{Yang:23arxiv}, HGRN and HGRN2 (both RNNs)~\citep{Qin:23,Qin:24arxiv}, RetNet (linear Transformer architecture)~\citep{Sun:23arxiv}, Hyena (linear Transformer)~\citep{Poli:23}, along with xLSTM[1:0] and xLSTM[7:1] variants discussed in Section~\ref{sec:xlstm_blocks_notation}. 

Mixed precision was used for training across models, although for RWKV-5, RWKV-6, GLA, and HGRN2, this approach did not incorporate PyTorch’s built-in automated mixed precision capabilities (for further discussion, see Appendix Section~\ref{sec:appGenTrainProc}). We organize the approaches into the following categories: (a) Transformers, (b) State Space Models (SSMs), and (c) Recurrent Neural Networks (RNNs) together with linear Transformers, where linear Transformers denote those architectures that implement attention using linear operations as a replacement for the standard Transformer attention mechanism. All architectures are designed to correspond to a 350M parameter GPT-3 reference, specifically with embedding dimension 1024 and 24 residual blocks. Of note, only GPT-3 employs weight sharing for token and output embeddings, leading to a reduced total parameter count for that architecture. Results summarized in Table~\ref{tab:spaj15b_model_results} demonstrate that xLSTM achieves the best validation perplexity among all tested methods. Additional information is available in Appendix~\ref{sec:appExpComparison}. Furthermore, Figure~\ref{fig:temp_sclaw15B} illustrates the scaling trends observed, which suggest that xLSTM is expected to maintain superior performance at larger model scales.

\paragraph{Ablation Studies.}
As illustrated in Table~\ref{tab:spaj15b_model_results} and Figure~\ref{fig:temp_sclaw15B}, xLSTM demonstrates strong language modeling capabilities when trained on 15B tokens derived from SlimPajama. To dissect the contributions of each architectural modification leading from LSTM to xLSTM, we sequentially transform a standard LSTM model into the full xLSTM design. The first modification incorporates LSTM layers within a pre-LayerNorm residual backbone. Next, this setup is expanded by introducing a post up-projection block. In the final step, we enhance the model with exponential gating mechanisms and matrix memory.

The results are shown in Table~\ref{tab:ablstudies} (top). 

Ablation experiments reveal that the notable gains in performance can be credited to both the exponential gating mechanism and the incorporation of matrix memory. Given the significance of gating structures in RNNs and State Space Models, we conduct ablations comparing alternative gating strategies. As reported in Table~\ref{tab:ablstudies} (bottom), our findings show that learning each gate parameter and allowing input-driven modulation consistently enhances results. Further analysis regarding the distinct contributions of backbone elements is provided in Appendix~\ref{sec:appAblDetails}.

\subsection{xLSTM as Large Language Model}
\label{sec:ExpLanguage}

To conclude our investigation, we conduct extensive language modeling experiments to assess xLSTM’s capabilities as a large language model. Accordingly, we scale up the dataset, training with 300B tokens sourced from SlimPajama. This token count matches those used in studies such as Mamba~\citep{Gu:24arxiv} and Griffin~\citep{De:24arxiv}.

Our evaluation positions xLSTM alongside RWKV-4, Llama, and Mamba, with each method chosen as a representative of its category described in Section~\ref{sec:ExpComparison}. RWKV-4 is included as the RNN baseline, given that suitable training precision parameters for RWKV-5, RWKV-6, and HGRN2 were only established after the commencement of 300B token training runs (as detailed in Appendix~\ref{sec:appGenTrainProc}). 

We experiment with various model scales (125M, 350M, 760M, and 1.3B), systematically examining each model's ability to generalize to longer sequences and measuring their effectiveness using the validation set. Additionally, we gauge their downstream task performance, analyze their language modeling competence across 571 textual domains featured in the PALOMA benchmark, and study the characteristics of their scaling laws.

\paragraph{Sequence Length Extrapolation.}
\label{sec:spaj300B_ctx_extrapolation}
To begin, we evaluate the sequence length extrapolation abilities of large models, specifically those with 1.3B parameters, including xLSTM, RWKV-4, Llama, and Mamba. These models are all trained with a context window of 2048 tokens, but are subsequently assessed on much larger context lengths, extending up to 16384. The resulting performance is depicted in Figure~\ref{fig:spaj300B_seq_extrapolation}. Unlike alternative approaches, xLSTM is able to sustain low perplexity even as the context length increases.

\paragraph{Validation Perplexity and Downstream Tasks.}
\label{sec:spaj_downstream_eval}
In addition, across each model scale, we assess xLSTM, RWKV-4, Llama, and Mamba models on next token prediction using the SlimPajama validation dataset, as well as on downstream benchmarks designed to evaluate capabilities in common sense reasoning.

The third column in Table~\ref{tab:lmeval} presents the perplexity scores on the validation set for various approaches. Among these, both xLSTM[1:0] and xLSTM[7:1] achieve the lowest validation perplexity across every evaluated model size. The additional columns in Table~\ref{tab:lmeval} display outcomes on several downstream tasks. In nearly every task and for all model capacities, xLSTM consistently outperforms other methods—Mamba surpasses xLSTM only on certain occasions within the ARC benchmark. Further information is available in Appendix~\ref{sec:appExpLanguage}.

\paragraph{Performance on PALOMA Language Tasks.} Third, we evaluate the next token prediction capabilities of xLSTM, RWKV-4, Llama, and Mamba models across all model sizes, specifically focusing on PALOMA language tasks~\citep{Magnusson:23arxivshort}. The assessment relies on perplexity scores computed for next token prediction spanning 571 diverse text domains, including sources such as nytimes.com and the Reddit community r/depression. Table~\ref{tab:ppleval} reports the grouped perplexity results, distinguishing between broader language modeling categories (first seven columns) and more fine-grained domain-specific benchmarks (last five columns). Notably, the xLSTM[1:0] variant yields superior results compared to xLSTM[7:1] on these tasks.

xLSTM[1:0] achieves lower perplexity compared to Mamba in 568 of 571 text domains (99.5\%), outperforms Llama in 486 out of 571 cases (85.1\%), and surpasses RWKV-4 in 570 out of 571 domains (99.8\%). Further information is available in Appendix~\ref{sec:appExpLanguage}.

\paragraph{Scaling Laws.} Fourth, we investigate the power-law scaling characteristics, which facilitate performance projection to models of greater scale~\citep{Kaplan:20,Brown:20short}.
In Figure~\ref{fig:temp_sclaw300B}, the scaling trends are displayed. While all the evaluated models exhibit comparable scaling trends, their performance differs by constant offsets. Among them, RWKV-4 lags the most, with Llama and Mamba performing moderately better. Notably, xLSTM surpasses Mamba, demonstrating a performance gap akin to that observed between Mamba and Llama. The observed scaling results suggest that, as models become larger, xLSTM maintains a consistent advantage relative to both Transformer-based and State-Space models.

\textbf{Generation Times and Maximal Throughput.} We evaluate the duration required for text generation in Figure~\ref{fig:inference_time_generation} and assess maximal throughput in Figure~\ref{fig:inference_time_decoding_throughput} (left) using our xLSTM models with 1.3B parameters. These results are juxtaposed with equivalent-sized implementations of Mamba, Llama, and RWKV provided by HuggingFace, where we also employ a static key-value cache for Llama. At the time these tests were conducted, attempts to utilize full cache compilation for the Transformer Model and to compile Mamba with \texttt{torch.compile} were unsuccessful. In generation time assessments, every model operates at batch size 1 with a pre-fill of 16 tokens, which should optimize performance for the Transformer architecture.

As illustrated in Figure~\ref{fig:inference_time_generation}, both xLSTM and the recurrent baselines Mamba and RWKV-4 exhibit linear growth with respect to sequence length, while Llama demonstrates quadratic scaling. Throughput under various batch sizes and pre-fill values is specifically evaluated for Llama, as detailed in Figure~\ref{fig:inference_time_decoding_throughput} (right). Here, the results indicate that xLSTM maintains substantially higher batch size capabilities compared to Llama—attributable to its constant memory profile—ultimately enabling xLSTM to reach superior throughput levels.

\section{Limitations}

(i) Unlike mLSTM, the architecture of sLSTM prevents the use of parallel memory mixing, rendering highly parallel implementations infeasible. Despite this, we engineered an efficient CUDA kernel tailored for sLSTM, achieving performance that is less than twice as slow as our optimized parallel mLSTM kernel. (ii) It should be noted that mLSTM's CUDA kernels have yet to be extensively optimized, making our current version approximately four times slower relative to FlashAttention or the scan technique integrated within Mamba. With further optimization, akin to approaches used in FlashAttention, significant speed gains are likely achievable. (iii) The computational demands of mLSTM’s matrix memory are considerable due to operations involving $d \times d$ matrices. However, since the procedures for updating and accessing memory are free from learnable parameters and amenable to conventional matrix parallelism, the actual impact on runtime remains modest. (iv) Careful attention must be paid to initializing the forget gates appropriately. (v) The capacity of the matrix memory is uncoupled from the sequence length, meaning that very long sequences may risk exceeding its storage limits. Empirically, this does not restrict operation on sequences up to 16k tokens; refer to Section~\ref{sec:spaj300B_ctx_extrapolation} for more details. (vi) Given the high computational costs associated with large-scale language model experiments, neither the model architecture nor the set of hyperparameters, particularly for larger xLSTM instances, underwent exhaustive optimization. We expect that thorough tuning will be essential for xLSTM to realize optimal performance.

\section{Conclusion}
Our investigation has provided a partial response to the central question: To what extent can LSTM models advance in language modeling when scaled to billions of parameters? Based on current findings, it is evident that: ``They can reach at least the performance levels of established architectures such as Transformers and State Space Models.'' By augmenting the standard LSTM to create xLSTM—introducing exponential gating with memory mixing alongside a novel memory design—we observe that xLSTM attains strong results on language modeling benchmarks in direct comparison with leading-edge approaches like Transformers and State Space Models. The observed scaling behavior suggests that increasingly large xLSTM variants could become prominent contenders against today’s dominant Large Language Models constructed with Transformer architectures. Furthermore, xLSTM holds significant promise for broad applicability across domains such as Reinforcement Learning, Time Series Forecasting, and the representation of physical systems.

\end{document}